\bookchapter{第四十三章：對話工程——同AI高效協作嘅對話技術}

\chapterintro{對話係 AI 編程嘅核心互動方式。但係留意：喺呢度「對話」唔等於「傾偈」——高效嘅 AI 編程對話係一種有結構嘅溝通技術，我哋叫佢做「對話工程」。佢研究嘅係：點樣用最少嘅對話輪次，攞到最高質素嘅代碼產出。講白啲，呢一章就係教你點樣同 AI「講清楚啲嘢」。}

\booksection{一、Prompt 分層策略}

同 AI 嘅對話可以分做三個層次，每層負責唔同嘅功能。明白呢三個層次，你就可以好似砌積木咁組合出高效嘅 Prompt。

\begin{center}
\begin{tabularx}{\linewidth}{|X|X|X|X|}
\hline
\textbf{層次} & \textbf{功能} & \textbf{示例} & \textbf{頻率} \tabularnewline
\hline
系統層 & 設定AI嘅角色同全局約束 & 「你係一個React專家，代碼必須用TypeScript」 & 項目開始時設定一次 \tabularnewline
任務層 & 描述而家要做嘅具體任務 & 「整一個表單組件，包含用戶名同密碼欄位」 & 每個任務開始時畀出嚟 \tabularnewline
回饋層 & 對AI嘅產出作出評價同修正 & 「密碼欄位缺少顯示/隱藏切換，請加上」 & 根據產出質素畀出嚟 \tabularnewline
\hline
\end{tabularx}
\end{center}

\begin{center}\small\color{gray}對話 Prompt 嘅三個層次\end{center}

系統層嘅 Prompt 通常寫喺項目記憶檔案入面（例如 CLAUDE.md），唔需要每次對話都重複。任務層係你對話時最常用嘅——佢應該包含「做乜嘢」「用乜嘢技術」「有乜嘢約束」三個要素。回饋層就最考功夫——好嘅回饋可以令 AI 精準修正，衰嘅回饋只會令佢重新估過。

\begin{tcbinfo}{小貼士：回饋三板斧}
一個實用嘅回饋模板係「指出問題 + 講明期望 + 畀上下文」。例如：「第 25 行嘅錯誤處理只係用咗 console.log（指出問題），生產環境應該用統一嘅錯誤處理器，並且傳回用戶友好提示（講明期望），項目入面有一個 handleError 函數可以用（畀上下文）」。呢種回饋比「改嚇錯誤處理啦」有效十倍。

\end{tcbinfo}

\booksection{二、除錯對話：點樣描述錯誤等 AI 快速定位}

除錯係 AI 編程入面最常見嘅場景之一。當你撞到 bug 嗰陣，點樣向 AI 描述問題，直接決定咗佢能否快速幫你定位同修復。以下係除錯對話嘅四個要點：

\begin{enumerate}
  \item 畀 AI「睇」原始資訊，而唔係你自己嘅轉述。將錯誤日誌、異常堆疊、相關代碼直接 @ 引用畀 AI，而唔係用「好似有個乜嘢報錯」嚟描述。AI 睇到原始資訊會比你嘅理解更準確。
  \item 描述「期望行為」同「實際行為」嘅分別。「我期望㩒咗按鈕之後會跳去首頁（期望），但實際㩒完個頁面都冇反應（實際）」——呢種描述比「個按鈕壞咗」有用得多。
  \item 提供重現步驟。「先輸入用戶名，然後㩒登入按鈕，就會出現呢個錯誤」——重現步驟等 AI 可以模擬你嘅操作路徑，更容易定位問題。
  \item 講清楚你已經試過乜嘢。「我已經檢查咗網絡請求，回傳 200，數據正常」——咁樣可以幫 AI 排除已經驗證過嘅方向，唔使嘥時間重複排查。
\end{enumerate}

差嘅除錯描述

\begin{lstlisting}
㩒個按鈕冇反應，幫我睇嚇
\end{lstlisting}

好嘅除錯描述

\begin{lstlisting}
㩒登入按鈕之後個頁面冇反應。 期望行為：㩒完之後跳去 /home 頁面 實際行為：㩒完完全冇反應，控制檯冇報錯 重現步驟： 1. 喺登入頁輸入用戶名admin同密碼123456 2. 㩒「登入」按鈕 3. 個頁面冇反應 已排查： - 網絡請求正常，回傳200 - onClick事件已綁定 - 檢查過LoginButton組件嘅代碼 相關檔案：@src/components/LoginButton.tsx @src/pages/Login.tsx
\end{lstlisting}

\booksection{三、迭代策略：小步快跑 vs 大步重構}

AI 編程入面嘅迭代有兩種基本策略，適用於唔同場景。明白佢哋嘅分別，可以幫你喺唔同情況下做出正確選擇。

「小步快跑」策略適用於功能開發同 bug 修復。佢嘅核心係：每次叫 AI 只做一個小修改，驗證通過先再做下一個。優勢係每一步都可以驗證、可以回退，風險可控；劣勢係速度慢啲，需要頻密嘅對話互動。適合：新功能開發、唔肯定邊個方案最好時嘅探索、對 AI 產出唔係好放心時嘅謹慎迭代。

「大步重構」策略適用於清楚知道目標嘅大規模修改。佢嘅核心係：一次過描述完整嘅目標，畀 AI 自主規劃同執行成個修改。優勢係速度快，AI 睇到全局圖景，可以做出更好嘅架構決策；劣勢係風險高，如果 AI 理解有偏差，修正成本好大。適合：明確嘅架構升級、大規模代碼遷移、有充分測試保障嘅重構。

\begin{center}
\begin{tabularx}{\linewidth}{|X|X|X|}
\hline
\textbf{維度} & \textbf{小步快跑} & \textbf{大步重構} \tabularnewline
\hline
每次修改量 & 一個功能點或者一個bug & 一個模組或者一個架構層 \tabularnewline
驗證頻率 & 每次修改完即刻驗證 & 全部修改完先統一驗證 \tabularnewline
對話輪次 & 多（5-10輪/功能） & 少（1-3輪/模組） \tabularnewline
風險 & 低，每步可以回退 & 高，糾偏成本大 \tabularnewline
適用場景 & 新功能、探索性開發 & 目標明確嘅重構、遷移 \tabularnewline
前提條件 & 冇特別要求 & 需要有測試覆蓋 \tabularnewline
\hline
\end{tabularx}
\end{center}

\begin{center}\small\color{gray}兩種迭代策略嘅對比\end{center}

實踐入面，兩種策略成日會夾埋用。例如重構一個模組嗰陣，先用「大步重構」叫 AI 畀個整體方案出嚟，審查完方案，再用「小步快跑」逐步執行。咁樣既保證咗全局視角，又控制咗執行風險。

\booksection{四、多輪對話嘅上下文管理}

隨住對話輪次增加，你需要主動管理對話上下文，避免「AI 忘記咗早期資訊」嘅問題。以下係幾個實用嘅上下文管理技巧。

第一，「錨定」關鍵資訊。對話開始嗰陣，將關鍵約束寫清楚，並且喺之後嘅對話中適時重申。例如「記住，我哋嘅數據存喺 localStorage，唔可以用數據庫」——喺相關嘅修改請求入面重申呢個約束，防止 AI 唔記得。

第二，善用「總結-續接」技巧。當對話變得好長嗰陣，可以叫 AI「總結嚇而家嘅進展同待辦事項」，然後開一個新對話，將呢個總結做起點。咁樣既保留咗關鍵上下文，又避免咗長對話嘅退化。

第三，「一個對話一個任務」。呢個係最簡單亦都最有效嘅上下文管理策略。每個對話只處理一個任務或者一個功能點，做完就開新對話處理下一個。咁樣自然限制咗對話長度，保持上下文清晰。

第四，善用工具嘅上下文管理功能。大部分 AI 編程工具都有「清除上下文」或者「新會話」掣——當你覺得 AI 嘅回答質素下降嗰陣，唔使猶豫，開一個新會話啦。唔見咗啲對話歷史，好過繼續喺退化咗嘅上下文入面掙紮好多。

\begin{tcbpractice}{嚟動手：做一次你 AI 嘅「病人」}
特登喺瀏覽器入面整一個小頁面，製造一個 bug（例如個掣冇反應）。然後跟本章「好嘅除錯描述」嘅格式，向 AI 完整描述一次：原始資訊、期望 vs 實際、重現步驟、已排查項。感受嚇 AI 定位問題嘅速度分別——你就會明白點解「描述問題」本身就係一門手藝。

\end{tcbpractice}

\begin{tcbsummary}{本章小結}
\begin{itemize}
  \item Prompt 三層：系統層（寫喺記憶檔案）、任務層（做乜嘢+技術+約束）、回饋層（指出問題+期望+上下文）
  \item 除錯四要點：畀原始資訊、期望vs實際、重現步驟、已排查項
  \item 小步快跑：低風險高輪次，適合新功能；大步重構：高風險低輪次，需要測試保障
  \item 上下文管理四招：錨定關鍵資訊、總結-續接、一個對話一個任務、果斷開新會話
\end{itemize}

\end{tcbsummary}
